% GENERATED FILE. Edit canonical lessons, then run npm run generate:books.
% canonical-source-hash: fnv1a64:e3c2968decd5d5ab
% canonical-lessons: JA-C489-nougyou, JA-C489-shuukaku, JA-C489-tagayasu, JA-C489-minoru, JA-C489-komugi

\chapter{Fields and Harvests}
\label{ch:ja-fields-and-harvests}
\hlchaptermodality{489}

\begin{chapteropening}
\textbf{By the end of this chapter:} \emph{I can talk about farming through the year.}

Complete the last lesson of chapter 489: I can talk about farming through the year.
\end{chapteropening}

\section[\texorpdfstring{\ja{のうぎょう} (nōgyō)}{のうぎょう (nōgyō)}]{\ja{のうぎょう} (nōgyō) — farming; agriculture}

\label{lesson:JA-C489-nougyou}



\begin{quote}
Before the new one: say the Japanese for iron, then the Japanese for construction work.
\end{quote}

\subsection*{You'll want to know: \ja{のうぎょう}}
\textbf{\ja{のうぎょう}} — \emph{nōgyō} — \textquotedblleft{}farming; agriculture\textquotedblright{}.

\begin{quote}
\textbf{\ja{そふは} \ja{のうぎょうを} \ja{しています。}} — \emph{sofu wa nōgyō o shiteimasu} — \textquotedblleft{}My grandfather is a farmer.\textquotedblright{}
\end{quote}

\begin{cousinweb}
\textbf{\ja{のうぎょう}} (\emph{nōgyō}) is Sino-Japanese (kango): two Chinese characters read with their on'yomi, readings borrowed from Chinese. They are \emph{nō}, \textquotedblleft{}agriculture\textquotedblright{} (Middle Chinese \emph{nowng}) and \emph{gyō}, \textquotedblleft{}business\textquotedblright{} (\emph{ngjaep}) (KANJIDIC2; Middle Chinese in Baxter's transcription).
\end{cousinweb}

\begin{grammarlens}[title={What you've built}]
One more word for describing what you lived and felt.
\end{grammarlens}

\subsection*{Your turn}
\begin{itemize}
\raggedright
  \item \emph{Say it:} \emph{nōgyō}
  \item \emph{Say it:} \emph{nōgyō}, once more
  \item \emph{Say it:} \emph{sofu wa nōgyō o shiteimasu} — My grandfather is a farmer.
  \item \emph{Recall:} say the Japanese for iron, then the Japanese for construction work, then say \emph{nōgyō} again
\end{itemize}

\begin{tcolorbox}[breakable,colback=teal!4,colframe=teal!35!black,title={Before you move on}]
What does \ja{のうぎょう} mean? (\textquotedblleft{}Farming; agriculture\textquotedblright{}.) Say it once more.
\end{tcolorbox}
\section[\texorpdfstring{\ja{しゅうかく} (shūkaku)}{しゅうかく (shūkaku)}]{\ja{しゅうかく} (shūkaku) — a harvest; a crop}

\label{lesson:JA-C489-shuukaku}



\begin{quote}
Before the new one: say the Japanese for construction work, then the Japanese for farming.
\end{quote}

\subsection*{You'll want to know: \ja{しゅうかく}}
\textbf{\ja{しゅうかく}} — \emph{shūkaku} — \textquotedblleft{}a harvest; a crop\textquotedblright{}.

\begin{quote}
\textbf{\ja{あきは} \ja{しゅうかくの} \ja{きせつです。}} — \emph{aki wa shūkaku no kisetsu desu} — \textquotedblleft{}Autumn is harvest season.\textquotedblright{}
\end{quote}

\begin{cousinweb}
\textbf{\ja{しゅうかく}} (\emph{shūkaku}) is Sino-Japanese (kango): two Chinese characters read with their on'yomi, readings borrowed from Chinese. They are \emph{shū}, \textquotedblleft{}income\textquotedblright{} (Middle Chinese \emph{syuw}) and \emph{kaku}, \textquotedblleft{}harvest\textquotedblright{} (\emph{hwak}) (KANJIDIC2; Middle Chinese in Baxter's transcription).
\end{cousinweb}

\begin{grammarlens}[title={What you've built}]
One more word for describing what you lived and felt.
\end{grammarlens}

\subsection*{Your turn}
\begin{itemize}
\raggedright
  \item \emph{Say it:} \emph{shūkaku}
  \item \emph{Say it:} \emph{shūkaku}, once more
  \item \emph{Say it:} \emph{aki wa shūkaku no kisetsu desu} — Autumn is harvest season.
  \item \emph{Recall:} say the Japanese for construction work, then the Japanese for farming, then say \emph{shūkaku} again
\end{itemize}

\begin{tcolorbox}[breakable,colback=teal!4,colframe=teal!35!black,title={Before you move on}]
What does \ja{しゅうかく} mean? (\textquotedblleft{}A harvest; a crop\textquotedblright{}.) Say it once more.
\end{tcolorbox}
\section[\texorpdfstring{\ja{たがやす} (tagayasu)}{たがやす (tagayasu)}]{\ja{たがやす} (tagayasu) — to plough; to till}

\label{lesson:JA-C489-tagayasu}



\begin{quote}
Before the new one: say the Japanese for farming, then the Japanese for a harvest.
\end{quote}

\subsection*{You'll want to know: \ja{たがやす}}
\textbf{\ja{たがやす}} — \emph{tagayasu} — \textquotedblleft{}to plough; to till\textquotedblright{}.

\begin{quote}
\textbf{\ja{のうかは} \ja{はるに} \ja{たんぼを} \ja{たがやします。}} — \emph{nōka wa haru ni tanbo o tagayashimasu} — \textquotedblleft{}Farmers plough the rice fields in spring.\textquotedblright{}
\end{quote}

Polite \textbf{\ja{たがやします}}, \emph{tagayashimasu}; a godan verb. Past (polite): \textbf{\ja{のうかが} \ja{たんぼを} \ja{たがやしました。}}, \emph{nōka ga tanbo o tagayashimashita}.

\begin{cousinweb}
\textbf{\ja{たがやす}} (\emph{tagayasu}) is a native Japanese word (wago), inherited rather than borrowed from Chinese; its kanji was attached later as a kun'yomi, and its older origin is not settled (JMdict, KANJIDIC2).
\end{cousinweb}

\begin{grammarlens}[title={What you've built}]
One more word for describing what you lived and felt.
\end{grammarlens}

\subsection*{Your turn}
\begin{itemize}
\raggedright
  \item \emph{Say it:} \emph{tagayasu}
  \item \emph{Say it:} \emph{tagayasu}, once more
  \item \emph{Say it:} \emph{nōka wa haru ni tanbo o tagayashimasu} — Farmers plough the rice fields in spring.
  \item \emph{Recall:} say the Japanese for farming, then the Japanese for a harvest, then say \emph{tagayasu} again
\end{itemize}

\begin{tcolorbox}[breakable,colback=teal!4,colframe=teal!35!black,title={Before you move on}]
What does \ja{たがやす} mean? (\textquotedblleft{}To plough; to till\textquotedblright{}.) Say it once more.
\end{tcolorbox}
\section[\texorpdfstring{\ja{みのる} (minoru)}{みのる (minoru)}]{\ja{みのる} (minoru) — to ripen; to bear fruit}

\label{lesson:JA-C489-minoru}



\begin{quote}
Before the new one: say the Japanese for a harvest, then the Japanese for to plough.
\end{quote}

\subsection*{You'll want to know: \ja{みのる}}
\textbf{\ja{みのる}} — \emph{minoru} — \textquotedblleft{}to ripen; to bear fruit\textquotedblright{}.

\begin{quote}
\textbf{\ja{あきに} \ja{いねが} \ja{みのります。}} — \emph{aki ni ine ga minorimasu} — \textquotedblleft{}The rice ripens in autumn.\textquotedblright{}
\end{quote}

Polite \textbf{\ja{みのります}}, \emph{minorimasu}; a godan verb. Past (polite): \textbf{\ja{ことしは} \ja{りんごが} \ja{みのりました。}}, \emph{kotoshi wa ringo ga minorimashita}.

\begin{cousinweb}
\textbf{\ja{みのる}} (\emph{minoru}) is a native Japanese word (wago), inherited rather than borrowed from Chinese; its kanji was attached later as a kun'yomi, and its older origin is not settled (JMdict, KANJIDIC2).
\end{cousinweb}

\begin{grammarlens}[title={What you've built}]
One more word for describing what you lived and felt.
\end{grammarlens}

\subsection*{Your turn}
\begin{itemize}
\raggedright
  \item \emph{Say it:} \emph{minoru}
  \item \emph{Say it:} \emph{minoru}, once more
  \item \emph{Say it:} \emph{aki ni ine ga minorimasu} — The rice ripens in autumn.
  \item \emph{Recall:} say the Japanese for a harvest, then the Japanese for to plough, then say \emph{minoru} again
\end{itemize}

\begin{tcolorbox}[breakable,colback=teal!4,colframe=teal!35!black,title={Before you move on}]
What does \ja{みのる} mean? (\textquotedblleft{}To ripen; to bear fruit\textquotedblright{}.) Say it once more.
\end{tcolorbox}
\section[\texorpdfstring{\ja{こむぎ} (komugi)}{こむぎ (komugi)}]{\ja{こむぎ} (komugi) — wheat}

\label{lesson:JA-C489-komugi}



\begin{quote}
Before the new one: say the Japanese for to plough, then the Japanese for to ripen.
\end{quote}

\subsection*{You'll want to know: \ja{こむぎ}}
\textbf{\ja{こむぎ}} — \emph{komugi} — \textquotedblleft{}wheat\textquotedblright{}.

\begin{quote}
\textbf{\ja{こむぎの} \ja{こなで} \ja{うどんを} \ja{つくります。}} — \emph{komugi no kona de udon o tsukurimasu} — \textquotedblleft{}We make udon from wheat flour.\textquotedblright{}
\end{quote}

\begin{cousinweb}
\textbf{\ja{こむぎ}} (\emph{komugi}) is a native Japanese compound (wago), inherited rather than borrowed from Chinese: built from \emph{ko}, \textquotedblleft{}small\textquotedblright{} and \emph{mugi}, \textquotedblleft{}wheat, barley\textquotedblright{} (JMdict, KANJIDIC2).
\end{cousinweb}

\begin{grammarlens}[title={What you've built}]
One more word for describing what you lived and felt.
\end{grammarlens}

\subsection*{Your turn}
\begin{itemize}
\raggedright
  \item \emph{Say it:} \emph{komugi}
  \item \emph{Say it:} \emph{komugi}, once more
  \item \emph{Say it:} \emph{komugi no kona de udon o tsukurimasu} — We make udon from wheat flour.
  \item \emph{Recall:} say the Japanese for to plough, then the Japanese for to ripen, then say \emph{komugi} again
\end{itemize}

\begin{tcolorbox}[breakable,colback=teal!4,colframe=teal!35!black,title={Before you move on}]
What does \ja{こむぎ} mean? (\textquotedblleft{}Wheat\textquotedblright{}.) Say it once more.
\end{tcolorbox}
